\chapter{Instalación y primer arranque}\label{cap:instalar}

\begin{objbox}
Instalarás REM de dos maneras (desde el código o con la aplicación de Mac), recorrerás el asistente de primera vez y conocerás dónde se guardan tus datos.
\end{objbox}

\section{Requisitos}
\begin{itemize}
\item Un Mac (la aplicación nativa requiere macOS 13 o posterior). El servidor funciona también en otros sistemas, pero algunas funciones (llavero, voz del sistema, icono de la barra de menús) son de macOS.
\item Node.js 22.13 o posterior (solo si arrancas desde el código).
\item Opcional: Ollama (modelo de IA local), whisper.cpp (dictado de notas de voz).
\end{itemize}

\section{Desde el código fuente}
Con el repositorio descargado:
\begin{codebox}
cd rem-negocio
npm start
\end{codebox}
La aplicación queda en \url{http://127.0.0.1:4400}. Si quieres otra carpeta de datos u otro puerto:
\begin{codebox}
REM_DATA_DIR=/ruta/a/mis-datos REM_PORT=4500 npm start
\end{codebox}
Por defecto, los datos se guardan en \ruta{~/Library/Application Support/REM Negocio}, en un fichero SQLite, con una copia automática diaria.

\begin{plainbox}{127.0.0.1}
Es la dirección que significa «este mismo ordenador». REM solo escucha ahí, de modo que ningún otro equipo de la red puede entrar.
\end{plainbox}

\section{Con la aplicación de Mac}
En la carpeta \ruta{mac/} hay una aplicación nativa (SwiftUI) que lleva el servidor dentro, lo arranca y lo vigila, y te guía paso a paso la primera vez. Se compila con \ruta{scripts/build-mac-app.sh}. La primera vez muestra un asistente de siete pasos.

\figura[0.36]{m-paso-0-welcome}{Asistente de la aplicación de Mac: bienvenida.}
La bienvenida resume tres ideas (todo se queda en el Mac, la IA propone y tú decides, puedes parar y retomar) y hace una \textbf{comprobación del equipo}: que el servidor local funciona, que hay espacio libre en disco y si el disco está cifrado con FileVault. Cada paso es breve y se puede dejar para más tarde: lo que has hecho se conserva.

\movil[0.30]{m-paso-1-business.png,m-paso-2-billing.png}{Pasos 1 y 2: datos del negocio (con datos de ejemplo) y cómo facturas.}{asist12}
\begin{description}[leftmargin=0.6cm,style=nextline]
\item[Tu negocio.] Nombre o razón social, NIF/CIF, dirección, correo y teléfono, y un logotipo opcional. Esos datos aparecen en tus facturas y presupuestos.
\item[Cómo facturas.] Eliges el tipo de actividad (servicios profesionales, oficios y reformas, comercio, hostelería u otro) y REM propone los valores habituales: IVA, retención de IRPF, plazo de pago, validez de los presupuestos, prefijo de las facturas, tarifa por hora e IBAN para los cobros. Todo se puede cambiar después.
\end{description}

\movil[0.30]{m-paso-3-backup.png,m-paso-4-data.png}{Pasos 3 y 4: copias de seguridad y traer tus datos.}{asist34}
\begin{description}[leftmargin=0.6cm,style=nextline]
\item[Copias de seguridad.] Eliges una carpeta (un disco externo, un NAS o iCloud Drive). Las copias se guardan \textbf{cifradas}, una cada día y otra al cerrar. Al final, el asistente \textbf{prueba una restauración}: una copia sin probar no es una garantía.
\item[Traer tus datos.] Tres vías: leer los \emph{Contactos de este Mac} (con vista previa de lo nuevo y lo repetido), abrir el \emph{importador} de archivos de otros programas (Google, Outlook, vCard, Excel, facturas emitidas y gastos) y, opcionalmente, usar tu propia clave de la API de Claude para que una IA proponga cómo traducir un archivo raro; antes de enviar nada se te muestra exactamente qué se enviaría, enmascarado.
\end{description}

\movil[0.30]{m-paso-5-access.png,m-paso-6-done.png}{Pasos 5 y 6: acceso y arranque, y fin del asistente.}{asist56}
\begin{description}[leftmargin=0.6cm,style=nextline]
\item[Acceso y arranque.] PIN opcional, arranque automático con el Mac e icono en la barra de menús.
\item[Todo listo.] Una lista de lo que queda pendiente, con un botón \ui{Completar} en cada punto, y un resumen de cómo se usa REM en cuatro líneas. El botón \ui{Abrir REM Negocio} termina el asistente.
\end{description}

\section{Primer arranque en el navegador}
Si arrancas sin datos, REM muestra una pantalla de bienvenida con tres pasos (nombre del negocio, valores habituales y cómo empezar). Puedes elegir \ui{Explorar con datos de ejemplo}, que carga un negocio ficticio para probar sin riesgo, y borrarlo después desde \ui{Ajustes → Datos y copias}.

\begin{checkbox}
Abre REM, elige \ui{Explorar con datos de ejemplo} y localiza en la barra lateral las secciones \ui{Hoy}, \ui{Mensajes}, \ui{Clientes} y \ui{Dinero}. Las verás con detalle en los capítulos siguientes.
\end{checkbox}
